\begin{lemma}[The replacement presentation carries the normalized datum]
Fix $C$, its minimum context, and raw Case Three data $D$. Suppose $\beta$
is missing on $B=T(y_{p-2})$, unused on the entire exact tail above
$A=T^*_{n,q}$, and satisfies the three-way alternative: when $q>0$ it is
outside $D_{n,q}\cup\Gamma^q$; or when $q=0$ it is outside the raw
terminal computed set $D_n$ and $\Gamma^0$; or it belongs to $D_{n,q}$
and its Gamma pair is disjoint from all full-tail edge colors.
Let $P$ be a replacement presentation on the same exact $T$. Assume
all displayed edges of $T$ retain their old values under $P$'s coloring,
the coloring is invariant on $B^+=T(y_{p-1})$, and $\beta$ is missing at
$y_p$ under it. Then there is a normalized Case Three datum for $P$
whose color is $\beta$.
\end{lemma}
\begin{proof}
By \noderef{HierarchyReplacementPresentation}, $P$ has the same-good
relation, a corresponding full-tree ETT and the transported minimum
context. Apply \noderef{P2752_T2_HierarchyMeasures} to the same-stage
part of \noderef{SameGoodHierarchy}. The old and new last starred stages
agree, as do their path numbers. The subtraction in
\noderef{HierarchyFullTailLength} consequently gives equal full-tail
lengths. For every numerical index $i$, the definition
\noderef{HierarchyFullTailPrefix} now uses the same terminal tree, the
same starred length and the same index, hence defines the same exact
prefix in both presentations.

Set $\ell$ equal to the common full-tail length. The common datum in
\noderef{HierarchySituationTwoCommonTailData} gives $\ell\ge2$.
The prefix formulas give $|B|=|T|-2$, $|B^+|=|T|-1$ and the exact
relations $B\preccurlyeq B^+\preccurlyeq T$. To verify the first
relation, compare at each retained index the two entries with the same
entry of $T$, as in \noderef{TreeSequenceInitialSegment}.
Every vertex of $B$ therefore belongs to $B^+$ by
\noderef{TreeSequenceInitialSegmentVertexSetSubset}. Expand the assumed
tree-missing membership using \noderef{TreeMissingColors} and take a
specific vertex $v\in V(B)$ at which $\beta$ is missing. The vertex
clause of \noderef{TreeColoringInvariant} on $B^+$ preserves its missing
set. Thus $\beta$ is missing at this same $v$ under $P$'s coloring and is
tree-missing on the common $B$.

By \noderef{TreeSequenceVertexSet}, take an exact displayed index of
$v$ in $B$ and cast it to its identical index in $T$ using the prefix
relation. Its value is at most $|B|=|T|-2$. The stored penultimate index
has value $|T|-1$, so this new earlier index is strictly before it.
Use it as the earlier witness in the new
\noderef{HierarchySituationTwoCaseThreeData}. Set the new alpha to
$\beta$. Earlier missingness is the just-preserved witness and terminal
missingness is the hypothesis. For all common-tail fields copy the old
last edge, penultimate and terminal indices, their position equations,
nonincidence, and theta. Fixed-edge equality preserves theta on that
exact last edge. The old full-tail inequality transports by length
equality; the path-number-full equation transports by both path-number
and full-tail equality. This verifies every inherited common-tail and
raw Case Three field.

For the extra fields of
\noderef{HierarchySituationTwoCaseThreeNormalizedData}, tree-missingness
two positions before the end was proved above, and it is on the required
new prefix since prefixes agree. For full-tail nonuse, the new starred
tree equals the old one. At every edge index above it, fixed-edge
equality rewrites the new edge value as its old value; apply old nonuse.

Finally we transport the three-way alternative exactly. The divider counts
are equal by same-good. At the last internal index, the cast of the old
last index is the new last index since their numerical values are their
respective equal counts. The stored-set clause of same-good therefore
identifies the two $D_{n,q}$ sets. For every color in that common set its
Gamma values agree. Taking the indexed union identifies $\Gamma^q$.
Apply \noderef{P2752_T2_TerminalStoredColors} to $P$'s corresponding
certificate to identify the two computed raw $D_n$ sets, including the
missing-set subtraction in \noderef{TashkinovStoredColors}. The exact
full-tail color sets agree by common starred-stage length and fixed-edge
equality at every index. Thus in the positive-divider alternative the
count inequality and both nonmembership clauses transport; in the
zero-divider alternative the count equation and raw-set and Gamma
nonmembership clauses transport. In the stored alternative, membership
puts $\beta$ in the common stored set, so its particular Gamma pair
agrees, and the equal full-tail color set transports the disjointness
clause. These are exactly the three alternatives required for the new
normalized datum. All its fields are now supplied and its alpha is the
chosen $\beta$ by definition.
\end{proof}
